% =============================================================================
%  Section — MCP-Assisted Evaluation: Claude Sonnet 4.6 (ByteBell Knowledge Graph)
%  Standalone section; drop into the full report alongside section2_mcp_runs.tex
% =============================================================================

\section{MCP-Assisted Evaluation --- Claude Sonnet~4.6}\label{sec:mcp-sonnet}

\subsection{Methodology}

Each of the 22 \texttt{astropy/astropy} tasks from the
\texttt{princeton-nlp/SWE-bench\_Verified} dataset was presented to
Claude Code (model: \texttt{claude-sonnet-4-6}) with access restricted
\emph{exclusively} to ByteBell MCP tools.
The model could not read local files, run \texttt{git} commands, call
any remote REST/GraphQL API, or use web search.
All codebase understanding had to be derived through the ByteBell
knowledge graph --- a pre-indexed, semantically tagged representation of
the repository.

\paragraph{Prompt structure.}
Every task was issued under a minimal prompt with two components:

\begin{enumerate}
  \item \textbf{Hard constraint:} ``ONLY use ByteBell MCP tools --- no
        local file reads, no git APIs, no GitNexus skills.''
  \item \textbf{Answer format:} ``Just write your answer, no format
        needed.  \texttt{\{\ answer:\ "answer"\ \}}\ --- that's it.''
\end{enumerate}

\noindent
The SWE-Bench question text was appended directly after these rules.
Claude Sonnet~4.6 then explored the codebase through MCP tool calls
(\texttt{smart\_search}, \texttt{keyword\_lookup},
\texttt{retrieve\_file}, \texttt{graph\_traverse}, etc.) and wrote
its answer to \texttt{answer.json}.
No explicit instruction to produce a diff was given; the model chose
its own output format (descriptions, code blocks, or unified diffs)
per task.

\paragraph{Judging.}
Answers were graded against the ground-truth patch using the same
three-dimensional rubric as all other runs (Table~\ref{tab:rubric-mcp-sonnet}),
scored by a separate Claude Sonnet~4.6 judge session with access to
\texttt{astropy\_tasks.json}.

\begin{table}[h]
\centering
\caption{Scoring rubric --- 8-point scale.}\label{tab:rubric-mcp-sonnet}
\begin{tabular}{lcp{8cm}}
\toprule
\textbf{Dimension} & \textbf{Max} & \textbf{Criterion} \\
\midrule
Root cause identification   & 3 & Correctly diagnoses \emph{why} the bug exists. \\
Correct file(s)             & 2 & Identifies the right file(s) to change. \\
Correct patch / code change & 3 & Proposed change matches ground truth or is functionally equivalent. \\
\bottomrule
\end{tabular}
\end{table}

\noindent
Grade tiers: \textbf{Exact}~(8/8), \textbf{Near-perfect}~(7/8),
\textbf{Partial}~(5--6/8), \textbf{Fail}~($\le$4/8).

% ---------------------------------------------------------------------------
\subsection{Results}

\subsubsection{Aggregate Performance}

\begin{table}[h]
\centering
\caption{Dimension-level scores across all 22 MCP Sonnet tasks.}\label{tab:dim-mcp-sonnet}
\begin{tabular}{lrrr}
\toprule
\textbf{Dimension} & \textbf{Score} & \textbf{Max} & \textbf{\%} \\
\midrule
Root cause identification   & 65 & 66 & 98.5\% \\
Correct file(s)             & 44 & 44 & 100\%  \\
Correct patch / code change & 49 & 66 & 74.2\% \\
\midrule
\textbf{Overall}            & \textbf{158} & \textbf{176} & \textbf{89.8\%} \\
\bottomrule
\end{tabular}
\end{table}

\begin{table}[h]
\centering
\caption{Grade distribution (MCP Sonnet).}\label{tab:grade-mcp-sonnet}
\begin{tabular}{lcl}
\toprule
\textbf{Grade} & \textbf{Count} & \textbf{Tasks} \\
\midrule
Exact (8/8)        &  7 & 1, 3, 4, 6, 9, 14, 16 \\
Near-perfect (7/8) & 13 & 2, 5, 8, 10, 11, 12, 13, 15, 17, 18, 20, 21, 22 \\
Partial (5--6/8)   &  2 & 7, 19 \\
Fail ($\le$4/8)    &  0 & --- \\
\bottomrule
\end{tabular}
\end{table}

\subsubsection{Per-Task Scores, Time \& Cost}

\begin{table}[h]
\centering
\small
\caption{Per-task MCP Sonnet results --- score, wall-clock time, and API cost.}\label{tab:per-task-mcp-sonnet}
\setlength{\tabcolsep}{4pt}
\begin{tabular}{clccccrr}
\toprule
\textbf{\#} & \textbf{Instance ID} & \textbf{RC} & \textbf{Files} & \textbf{Patch} & \textbf{Score} & \textbf{Time (s)} & \textbf{Cost (\$)} \\
\midrule
 1 & \texttt{astropy-12907} & 3 & 2 & 3 & 8/8 &   369 & 0.996 \\
 2 & \texttt{astropy-13033} & 3 & 2 & 2 & 7/8 &   331 & 0.497 \\
 3 & \texttt{astropy-13236} & 3 & 2 & 3 & 8/8 &   355 & 0.584 \\
 4 & \texttt{astropy-13398} & 3 & 2 & 3 & 8/8 &    51 & 0.128 \\
 5 & \texttt{astropy-13453} & 3 & 2 & 2 & 7/8 &   121 & 0.357 \\
 6 & \texttt{astropy-13579} & 3 & 2 & 3 & 8/8 &    89 & 0.203 \\
 7 & \texttt{astropy-13977} & 3 & 2 & 1 & 6/8 &    66 & 0.147 \\
 8 & \texttt{astropy-14096} & 3 & 2 & 2 & 7/8 &    56 & 0.156 \\
 9 & \texttt{astropy-14182} & 3 & 2 & 3 & 8/8 &   142 & 0.299 \\
10 & \texttt{astropy-14309} & 3 & 2 & 2 & 7/8 &    45 & 0.121 \\
11 & \texttt{astropy-14365} & 3 & 2 & 2 & 7/8 &    80 & 0.158 \\
12 & \texttt{astropy-14369} & 3 & 2 & 2 & 7/8 &    80 & 0.203 \\
13 & \texttt{astropy-14508} & 3 & 2 & 2 & 7/8 &   202 & 0.358 \\
14 & \texttt{astropy-14539} & 3 & 2 & 3 & 8/8 &   269 & 0.432 \\
15 & \texttt{astropy-14598} & 3 & 2 & 2 & 7/8 &   385 & 0.575 \\
16 & \texttt{astropy-14995} & 3 & 2 & 3 & 8/8 &   100 & 0.254 \\
17 & \texttt{astropy-7166}  & 3 & 2 & 2 & 7/8 &    63 & 0.154 \\
18 & \texttt{astropy-7336}  & 3 & 2 & 2 & 7/8 &    89 & 0.211 \\
19 & \texttt{astropy-7606}  & 2 & 2 & 1 & 5/8 &    65 & 0.168 \\
20 & \texttt{astropy-7671}  & 3 & 2 & 2 & 7/8 &    41 & 0.085 \\
21 & \texttt{astropy-8707}  & 3 & 2 & 2 & 7/8 &   172 & 0.258 \\
22 & \texttt{astropy-8872}  & 3 & 2 & 2 & 7/8 &   103 & 0.255 \\
\midrule
   & \textbf{Total}  & \textbf{65/66} & \textbf{44/44} & \textbf{49/66} & \textbf{158/176} & \textbf{3{,}273} & \textbf{6.60} \\
   & \textbf{Average} &       &       &       & \textbf{7.2/8}   & \textbf{149} & \textbf{0.30} \\
\bottomrule
\end{tabular}
\end{table}

% ---------------------------------------------------------------------------
\subsubsection{Token \& Cost Breakdown (Aggregated)}

\begin{table}[h]
\centering
\small
\caption{Aggregated token usage and cost across all 22 MCP Sonnet tasks.}\label{tab:tokens-agg-mcp-sonnet}
\setlength{\tabcolsep}{4pt}
\begin{tabular}{lrrrrrrr}
\toprule
& \textbf{Input} & \textbf{Cache Write} & \textbf{Cache Read} & \textbf{Output}
& \textbf{Eff Weighted} & \textbf{Eff Input} & \textbf{Cost (\$)} \\
\midrule
Total   &        298 & 561{,}515 & 4{,}998{,}872 & 199{,}447 & 1{,}202{,}079 & 5{,}560{,}685 & 6.60 \\
Average &         14 &  25{,}523 &   227{,}221   &   9{,}066 &    54{,}640   &   252{,}759   & 0.30 \\
\bottomrule
\end{tabular}
\\[4pt]
{\footnotesize
Effective Weighted Input = Input + $1.25 \times$ Cache Write + $0.1 \times$ Cache Read.\quad
Effective Input = Input + Cache Write + Cache Read.}
\end{table}

% ---------------------------------------------------------------------------
\subsubsection{Per-Model Token Usage}

Unlike all other evaluated configurations --- Opus MCP and all raw runs ---
the Sonnet~4.6 MCP run dispatches \emph{all} requests to a single model.
There is no Haiku routing layer.
Every tool call, exploration step, and patch generation goes directly
through \texttt{claude-sonnet-4-6}.

\begin{table}[h]
\centering
\caption{Token and cost: Sonnet-only (no Haiku split).}\label{tab:model-split-mcp-sonnet}
\begin{tabular}{lrr}
\toprule
\textbf{Metric} & \textbf{Sonnet} & \textbf{Total} \\
\midrule
Input tokens       &         298 &         298 \\
Cache Write tokens & 561{,}515   & 561{,}515   \\
Cache Read tokens  & 4{,}998{,}872 & 4{,}998{,}872 \\
Output tokens      & 199{,}447   & 199{,}447   \\
Cost (USD)         & \$6.598     & \$6.598     \\
API requests       &         200 &         200 \\
\bottomrule
\end{tabular}
\end{table}

\noindent
All 200 API requests were made to Sonnet.  Average requests per task
is 9.1, versus 13.1 for Opus MCP (288 requests $\div$ 22 tasks).
The absence of a Haiku triage layer means Sonnet handles both the
lightweight routing decisions and the heavy reasoning within the same
model, but the MCP graph's targeted retrieval keeps total requests low.

% ---------------------------------------------------------------------------
\subsection{Observations}\label{sec:mcp-sonnet-obs}

The following observations are derived strictly from the Sonnet~4.6 MCP
run results and, where noted, from direct comparison with the Opus~4.6
MCP run (Section~\ref{sec:mcp}).

\begin{enumerate}

\item \textbf{Strong accuracy with zero filesystem access.}
      The model scored 158/176 (89.8\%) using only the ByteBell
      knowledge graph for codebase understanding.
      7 of 22 tasks received a perfect 8/8, and no task scored below
      5/8 --- preserving the zero-failure property seen in the Opus MCP
      run and improving on the raw Opus run which had one 3/8 fail.

\item \textbf{Root cause identification is the highest of all runs.}
      Root cause scored 65/66 (98.5\%), the best across all four
      evaluated configurations.  The single deduction came from
      task~19 (\texttt{7606}): Sonnet identified the \texttt{TypeError}
      symptom in \texttt{UnrecognizedUnit.\_\_eq\_\_} but missed the
      semantic requirement to return \texttt{NotImplemented} (not
      \texttt{False}) and the second required location in
      \texttt{UnitBase.\_\_eq\_\_}.

\item \textbf{File identification remains perfect.}
      All 22 tasks scored 44/44 on file identification.
      ByteBell's semantic indexing surfaces correct file paths
      without manual directory traversal, consistent with all other
      MCP and raw runs.

\item \textbf{Patch quality is the main weakness.}
      The patch dimension scored 49/66 (74.2\%), below the Opus MCP
      run (55/66, 83.3\%) and the Opus raw run (54/66, 81.8\%).
      Sonnet consistently produces functionally correct but
      non-minimal solutions: it over-engineers patches
      (\texttt{14508}, \texttt{14096}), uses equivalent but
      divergent API idioms (\texttt{7166}, \texttt{8872}), or takes
      alternative grammar/logic routes (\texttt{14369},
      \texttt{14598}).  Only tasks~7 (\texttt{13977}) and~19
      (\texttt{7606}) contain substantive correctness gaps.

\item \textbf{Sonnet MCP is 43\% cheaper per task than Opus MCP.}
      Average cost per task is \$0.30 (Sonnet MCP) versus \$0.52
      (Opus MCP), a 42\% reduction, and \$0.73 (Opus raw), a 59\%
      reduction.  Total cost across 22 tasks is \$6.60 (Sonnet MCP)
      versus \$11.53 (Opus MCP) and \$16.16 (Opus raw).

\item \textbf{No Haiku routing layer --- single-model execution.}
      All 200 requests went to Sonnet.  In both MCP and raw Opus
      runs, Haiku handled per-task routing stubs (22--178 requests
      across runs).  The Sonnet MCP harness does not instantiate a
      Haiku layer, confirming the cost and token savings are
      attributable to model tier, not to a routing change.

\item \textbf{Effective weighted input is 30\% lower than Opus MCP.}
      Total effective weighted input is 1.20\,M (Sonnet MCP) versus
      1.72\,M (Opus MCP) and 2.87\,M (Opus raw).  Sonnet's
      shorter-form answers (targeted code blocks rather than full
      unified diffs) are the primary driver: output tokens average
      9{,}066 per task (Sonnet MCP) versus 5{,}730 (Opus MCP).
      Cache write tokens average 25{,}523 (Sonnet MCP) versus
      32{,}000 (Opus MCP), reflecting fewer exploratory tool calls
      per task.

\item \textbf{Cache reads are 37\% lower than Opus MCP.}
      Total cache reads are 5.00\,M (Sonnet MCP) versus 7.87\,M
      (Opus MCP) and 14.79\,M (Opus raw).
      Sonnet achieves the same file identification accuracy as Opus
      with substantially fewer iterative reads, suggesting more
      decisive MCP query patterns.

\item \textbf{Wall-clock time is 40\% faster than Opus MCP.}
      Average time is 149\,s (Sonnet MCP) versus 249\,s (Opus MCP)
      and 251\,s (Opus raw).
      Unlike the Opus MCP vs.\ raw comparison where time was
      effectively identical, Sonnet's reduced output verbosity
      and fewer cache reads produce a measurable wall-clock
      advantage.  The maximum task time is 385\,s (\texttt{14598}),
      versus 845\,s (\texttt{14369}) in the Opus MCP run.

\item \textbf{The hardest task (13398) is solved in 51\,s at
      \$0.128 --- the fastest and cheapest result across all runs.}
      Task~4 (\texttt{astropy-13398}), rated 1--4\,h human
      difficulty, required adding topocentric ITRS transforms across
      six files including a new module.  Sonnet MCP returned a
      complete unified diff covering all six files in one pass,
      scoring 8/8.  For comparison: Opus MCP took 281\,s at \$0.70;
      Opus raw took 612\,s at \$1.77.  The MCP knowledge graph
      allowed Sonnet to locate all relevant coordinate-transform
      files in a single traversal rather than iteratively scanning
      the repository.

\item \textbf{No outright failures.}
      The Sonnet MCP run, like the Opus MCP run, has no task below
      5/8.  The Opus raw run had one 3/8 fail on task~19
      (\texttt{7606}) that produced no patch; Sonnet MCP scores
      5/8 on the same task, returning a partial fix with incorrect
      return semantics (\texttt{False} instead of
      \texttt{NotImplemented}) and missing the second location.

\item \textbf{Shared weakness: task~7 (\texttt{13977}).}
      Both Opus and Sonnet --- in every run mode --- score~6/8 on
      \texttt{Quantity.\_\_array\_ufunc\_\_} duck-type handling.
      Both wrap only the input-conversion loop in
      \texttt{try/except}, missing the wider scope required by the
      ground truth.  This is a consistent failure mode independent
      of model, context source, or cost tier, suggesting a genuine
      ambiguity in the issue statement rather than a knowledge-graph
      retrieval gap.

\item \textbf{Near-perfect dominates the grade distribution.}
      13 of 22 tasks (59\%) are Near-perfect (7/8), versus 6 of 22
      (27\%) in the Opus MCP run and 5 of 22 (23\%) in the Opus raw
      run.  Exact scores drop from 13--14 in Opus runs to 7 in
      Sonnet MCP.  This pattern is consistent with Sonnet finding
      functionally correct, non-minimal solutions rather than the
      exact idiomatic change present in the ground truth.

\end{enumerate}

% ---------------------------------------------------------------------------
\subsection{Cross-Run Summary}\label{sec:mcp-sonnet-crossrun}

Table~\ref{tab:crossrun} consolidates all four evaluated configurations
for direct comparison.

\begin{table}[h]
\centering
\small
\caption{Cross-run comparison across all four configurations.}\label{tab:crossrun}
\setlength{\tabcolsep}{5pt}
\begin{tabular}{lcccccrr}
\toprule
\textbf{Configuration} & \textbf{Score} & \textbf{\%} & \textbf{Exact} & \textbf{Near-perf} & \textbf{Fail} & \textbf{Avg Time (s)} & \textbf{Avg Cost (\$)} \\
\midrule
Opus raw          & 162/176 & 92.0\% & 14 & 5  & 1 & 251 & 0.734 \\
Opus MCP          & 163/176 & 92.6\% & 13 & 6  & 0 & 249 & 0.524 \\
Sonnet MCP        & 158/176 & 89.8\% &  7 & 13 & 0 & 149 & \textbf{0.300} \\
\bottomrule
\end{tabular}
\end{table}

\noindent
Sonnet~4.6 with MCP delivers the lowest cost (\$0.30/task) and fastest
wall-clock time (149\,s/task) of any configuration, while maintaining
competitive accuracy (89.8\%) and eliminating complete failures.
The trade-off is a shift from Exact to Near-perfect solutions: Sonnet
reliably finds the right files and root causes but tends toward
over-engineered or idiomatically divergent patches rather than the
minimal ground-truth change.
